\section{Discussion}
\label{sec:discussion}

\subsection{Conceptual implication}

Stability claims for local learners are incomplete unless they specify both the
perturbation operation and the evaluation point.  This paper has shown that
leave-one-out stability and replacement robustness, two notions that are often
used interchangeably in practice, are structurally distinct for deterministic
$k$-NN.  The distinction is not an asymptotic artifact: it appears in the
smallest possible finite metric configuration.

The core message is that LOO is a coupled deletion diagnostic, not a
certificate of replacement robustness.  A classifier can be perfectly
LOO-stable while being maximally vulnerable to a single adversarially chosen
replacement at a strategically chosen query.

\subsection{Practical implication}

Low LOO error should not be interpreted as evidence of replacement robustness.
For $k$-NN and related local rules, margin distributions are more informative
than aggregate LOO scores when assessing vulnerability to training-set
perturbations.  Queries with $|\M_k(S,x)| \le 2$ are potentially vulnerable to
a single label-flip replacement, and practitioners should monitor the margin
distribution rather than relying solely on LOO error.

\subsection{Limitations}

\begin{enumerate}[leftmargin=*,label=(\arabic*)]
    \item Binary classification with 0--1 loss.  Extending to real-valued
        scores, multiclass labels, or regression loss requires a more general
        notion of margin crossing.
    \item Finite metric spaces (often graph metrics).  Generalizing to infinite
        spaces, continuous densities, or non-metric similarity is open.
    \item Deterministic tie-breaking.  The even-$k$ results show that
        randomization can close some gaps; randomized tie-breaking is not
        analyzed.
    \item Worst-case rather than distributional analysis.  The paper studies
        fixed worst-case samples; connecting to expected or typical-case
        behavior is an open direction.
    \item The DLVR extension is not fully characterized.  Which local rules
        admit the separation is an open problem.
\end{enumerate}

\subsection{Open problems}

\begin{enumerate}[leftmargin=*,label=(\arabic*)]
    \item \textbf{Tie-free even-$k$ separation.}  Does there exist, for any
        even $k \ge 2$, a finite metric sample in which the original prediction
        at the designated query is determined by a strict majority (margin $\neq
        0$) and the LOO/replace-one separation holds?
    \item \textbf{Quantitative bounds.}  For a given data distribution, what is
        the expected LOO/replace-one gap?  Can the worst-case gap be bounded by
        a function of the data geometry?
    \item \textbf{Margin-based mitigation.}  Can the margin-based diagnostic
        be extended to an algorithm that repairs vulnerable queries, for
        instance by recommending a different $k$ or a specific subset of points
        to re-label?
    \item \textbf{Characterization of local vote rules.}  Which deterministic
        local vote rules admit an analogous LOO/replace-one separation?
\end{enumerate}
